\cotefile{157-1}
\scanpage{182}\mspage{38}
% archive sheet number: 182

Ainsi, les foncteurs familiers, \deleted{pour} une \deleted{flèche}

\[
u:I'\unsure{\longrightarrow|\mathbin{-}}I
\]

dans $\mathrm{Cat}$

\diagram{$(\unsure{3.7.})\quad\left\{\vphantom{\rule{0pt}{4em}}\right.$\begin{tikzcd}[ampersand replacement=\&,column sep=large]
I'{}^{\wedge} \arrow[r,shift left=4ex,"{u_{!}}"] \arrow[r,shift right=4ex,"{u_*}"] \& I^{\wedge} \arrow[l,"{u^*}"']
\end{tikzcd}$,$}

s'interprètent comme les foncteurs de \unsure{même nom|mêmes noms} pour le dérivateur $\mathrm{HOT}_{W_0}$, ou \inserted{plus} directement, pour le dérivateur (qui pouvait \unsure{justement|faussement|jusqu'ici} \unsure{être} considéré comme \unsure{canularesque}) $\mathbb{D}_{(\mathrm{Ens})}$.

Cela donne un sens à \deleted{\unsure{une autre}} la cohorte des notions cohomologiques associées à un dérivateur, ou plus particulièrement à un localiseur fondamental $W$ sur $\mathrm{Cat}$ (donnant naissance au dérivateur $\mathrm{HOT}_W$) \unsure{;} \unsure{On|les} \deleted{flèches} \deleted{$W_0$-équivalences \unsure{:}} les \unsure{désignera|désigner} par les noms de $W_0$-\unsure{foncté} (au lieu de $\mathrm{HOT}_{W_0}$-foncté, ou $\mathbb{D}_{\mathrm{Ens}}$-foncté), \deleted{\unsure{$\mathbb{D}$}} \deleted{\unsure{tels}} notamment pour une flèche $u:I'\longrightarrow I$ de $\mathrm{Cat}$ telle la notion

\[
\left\{\begin{aligned}
& W_0\text{-équivalence}\quad(\Longleftrightarrow\pi_0(u)\text{ bijectif})\\
& \vert\ \ W_0\text{-asphérique}\quad(\Longleftrightarrow\text{ les }I'_{\deleted{\unsure{u}}/i}\text{ pour }i\in\mathrm{Ob}\,I\text{ sont}\\[-2pt]
&\hspace{13em}0\text{-connexes}\\
& \vert\ \ W_0\text{-coasphérique}\quad(\Longleftrightarrow\text{ les }{}_{i\backslash}I'\text{ pour }i\in\mathrm{Ob}\,I\text{ sont}\\[-2pt]
&\hspace{13em}0\text{-connexes})\\
& W_0\text{-propre}\quad(\Longleftrightarrow\text{les inclusions }I'_i\longrightarrow I'_{/i}\ (\unsure{i'|i}\in\mathrm{Ob}\,\unsure{I|I'})\\
&\hspace{7em}\text{sont }W_0\text{-coasphériques, i.e. les fibres}\\
&\hspace{7em}\text{des }\deleted{\unsure{l}}\text{colocalisés }{}_{i'\backslash}I'_{\deleted{\unsure{/i}}}\longrightarrow{}_{i\backslash}I_{\deleted{\unsure{/i}}}\ (\unsure{i'=ui|i=ui'})\\
&\hspace{7em}\text{sont }0\text{-connexes}
\end{aligned}\right. \tag*{(3.8)}
\]

\margin{Ex : les Cat-\deleted{\unsure{fibrations}} cofibrations,}
\scanpage{183}\mspage{39}
% archive sheet number: 183

\margin{Ex : les Cat-fibrations}

$(3.8.)\quad$\raisebox{-5.5\baselineskip}[0pt][0pt]{$\left\{\vphantom{\rule[-5.5\baselineskip]{0pt}{11\baselineskip}}\right.$}\,$W_0$-lisse ($\Longleftrightarrow$ les inclusions $I'_i\unsure{\longrightarrow|\mathbin{-}}{}_{i\backslash}I'$ sont $W_0$-asphériques, i.e. les fibres des flèches localisées \deleted{$\unsure{A}$} $I'_{/i'}\unsure{\longrightarrow|\mathbin{-}}I_{/i}$ $(i=ui')$ sont $0$-connexes

$W_0$-parfait $=W_0$-propre et $W_0$-lisse

$W_0$-équivalences universelles ($\pi_0$ bijectif et le reste pour tout changement de base.)

\margin{Ex : les \deleted{\unsure{flèches}} $W_0$-lisses ou $W_0$-propres, à fibres $0$-connexes}

$\Downarrow$ impl. \uuline{stricte} (et comment)

\inserted{Ce qui impli\unsure{que} les \deleted{$W$}-fibrations (au \unsure{sens} faible) à fibres asphériques i.e. $0$-connexes}

$W_0$-fibration : $u^*:\mathrm{Cat}_{/I}\longrightarrow\mathrm{Cat}_{/I'}$ transforme $W_0$-équivalences en $W_0$-équivalences

\margin{\deleted{Ex : les \unsure{fibrations} \unsure{propres}}}

$\Downarrow$ impl. \uuline{stricte} (et comment !)

\inserted{$\longleftrightarrow$ mais si les fibres sont $0$-connexes, \unsure{équivalentes} et signifie \unsure{les changements|les équivalences} de \unsure{base} des $W_0$-\unsure{fibrations} sont des $W_0$-\unsure{équiv.} universelles\ldots}

\uline{$W_0$-fibrat\deleted{\unsure{ion}} faible} ($\Longleftrightarrow$ \deleted{pour tout} $u^*:\mathrm{Cat}_{/I}\longrightarrow\mathrm{Cat}_{/I'}$ transforme homotopismes $J_1\overset{\deleted{\unsure{\sim}}}{\longrightarrow}J_2$ en $W_0$-équivalences -- et il suffit de l'exiger pour l'inclusion d'un objet final \inserted{(ou, au choix, initial)} \deleted{\unsure{ou pour l'in}} et pour l'inclusion d'un point dans $\Delta^1$)

\margin{\uline{Ex} les fl. $W_0$-\deleted{\unsure{lisses}} telles que $\pi_0(I'/I)=(i\mapsto\pi_0({}_{i\backslash}I'))=u_!(e_{I'{}^\wedge})$ soit loc. constant sur $I$ i.e. $\underline{\pi}_0(I'/I)$ soit un \uuline{rev.} étale de $I$, ou si $u^{\mathrm{opp}}$ satisfait à cette condition}

Signalons, comme \inserted{\unsure{candide|corollaire}} utilisation de ces notions, que pour un diagramme cartésien dans $\mathrm{Cat}$

\diagram{$(3.9.)\qquad$\begin{tikzcd}[ampersand replacement=\&]
I' \arrow[d,"u"'] \& J' \arrow[l,"{p'}"'] \arrow[d,"v"]\\
I \& J \arrow[l,"p"]
\end{tikzcd}$,$}

la flèche de chgt de base pour $u_*,v_*$ \inserted{pour les préfaisceaux} est \unsure{is} si $u$ est \inserted{$W_0$-}propre, ou $p$ $W_0$-lisse, et

la flèche de \unsure{chgt.} de base pour $u_!,v_!$ est \unsure{id|iso} si $u$ est $W_0$-\deleted{\unsure{faible}}\unsure{lisse} ou $p$ $W_0$-propre.

% slanted
\margin{Sauf erreur, \unsure{cela se démontre|cela découle} \uuline{\unsure{sans|d'une}} \uline{\unsure{caractérisation|considération}} \unsure{très} \unsure{simple} des morphismes $W_0$-propres, $W_0$-lisses \unsure{uniquement|impliquant}}

Je signale à ce propos que ces énoncés, qui à priori se \deleted{\unsure{démontraient}} formulent pour les préfaisceaux à valeurs dans $\mathrm{Ens}$, \unsure{restent|restent-ils} valables
\scanpage{184}\mspage{40}
% archive sheet number: 184

pour des préfaisceaux à valeurs dans une catégorie $M$ quelconque (sous réserve d'existence des opérations \deleted{\unsure{de quelque nature}} d'image directe \inserted{cohomologique ou homologique} \unsure{envisagées|considérées}).

Cela se voit p. ex. en notant que

\[
W_0\text{-équivalence}\ \Longrightarrow\ \mathbb{D}_M\text{-équivalence}\tag*{(3.10.)}
\]

% slanted
\margin{\unsure{pour être tout|pour être} à fait \unsure{correct}, \inserted{\unsure{pour|parmi}} \deleted{il faut} ``\uline{\unsure{toute}}'' $W_0$-notion, excepter les deux notions de $W_0$-fibration (forte, et faible).}

pour tte \inserted{$\mathcal{U}$-}catégorie $M$, d'où résulte que \unsure{une|toute} \inserted{(*)} ``$W_0$-notion'' implique la $\mathbb{D}_M$-notion correspondante, et permet donc d'appliquer à $\mathbb{D}_M$ les résultats \unsure{le|la} concernant. (Il y a bien sûr un argument \unsure{direct moins|direct mais|d'autant} \unsure{savant|suivant}, consistant à se ramener au cas où $M$ est petite, et à utiliser alors les plongements pl. fidèles $M\hookrightarrow M^\wedge$ resp. $M\hookrightarrow((M^\circ)^\wedge)^\circ$.)

Je rappelle pour mémoire une (\unsure{troisième} ?) caractérisation des flèches $u:\deleted{\unsure{I}}I'\longrightarrow I$ $W_0$-propres (resp $W_0$-lisses) : ce sont celles \deleted{\unsure{p}} pour lesquelles le foncteur chgt. de base $u^*:\mathrm{Cat}_{/I}\longrightarrow\mathrm{Cat}_{/I'}$ transforme flèches $W_0$-coasphériques (resp. $W_0$-asphériques) en itou. Cela gagne en intérêt, du fait que la notion de foncteur $W_0$-coasphérique (resp. $W_0$-asphérique) s'interprète comme celle de foncteur cofinal (resp. co-cofinal), cf plus bas.
\scanpage{185}\mspage{41}
% archive sheet number: 185

La notion de $W_0$-fibration apparaît comme extrêmement forte -- beaucoup plus forte à certains égards que celle (\unsure{disons|absolue}) de $W_\infty$-fibration \inserted{(\unsure{du moins|au moins} \unsure{ostensiblement})} \deleted{$W_0$-}fibration. Elle \unsure{équivaut|revient} à ceci : pour tte composante connexe $I'_\alpha$ de $I'_\alpha$, désignant par $I_\alpha$ la composante connexe correspondante de $I$, et pour tout $J_\alpha$ $0$-connexe \inserted{et} au dessus de $I_\alpha$, le produit fibré $I'_\alpha\times_{I_\alpha}J_\alpha$ est encore $0$-connexe. (\unsure{Ça|Sa} ressemble aux histoires de totale $0$-connexité d'un topos !) \deleted{Cel\unsure{a} implique que les fibres de $I'_\alpha\longrightarrow I_\alpha$ sont $0$-connexes (prendre $J_\alpha=e_{\unsure{i}}$), et \unsure{a}} Mais cela signifie \deleted{\unsure{que|pour}} que $I'_\alpha\longrightarrow I_\alpha$ est \inserted{une $W_0$-équivalence} universelle\deleted{\unsure{ment}}. Cela implique que ses fibres sont $0$-connexes, et la réciproque est vraie si $u:I'\longrightarrow I$ est $W_0$-propre ou $W_0$-lisse. Ainsi, la condition \unsure{essentielle|en italiques}$\unsure{\Rightarrow|,}$ pour une \inserted{$W_0$-fibration} est celle que les fibres des flèches induites $I'_\alpha\longrightarrow I_{\unsure{\beta(\alpha)|\mu(\alpha)}}$ soient $0$-connexes ; \deleted{\unsure{si on suppose} \unsure{2\textsuperscript{o}} que pour une flèche de $I$ i.e.} \inserted{En fait}

\uuline{\deleted{\unsure{proposition}} Proposition 3.1} Soit $u:I'\longrightarrow I$ une flèche dans $\mathrm{Cat}$.

a) Pour que $u$ soit une $W_0$-équivalence universelle, il f et s. que \deleted{ses fibres soient $W_0$-connexes,} pour toute \deleted{flèche $\alpha$ de $I$} $\Delta^1\longrightarrow I$ i.e. \inserted{toute flèche $\alpha$ de $I$}, \deleted{et il suffit} la catégorie $I'\times_I\Delta^1\overset{\mathrm{déf}}{=}\deleted{\unsure{I}}I'|\alpha$ soit \deleted{\unsure{les points} $\alpha$ \unsure{mis} \unsure{bout}} $0$-connexe, ou encore que les fibres de $I$ soient $0$-connexes, et que pour toute flèche \uline{non identique} $\alpha$ de $I$, $I'|\alpha$ soit $0$-connexe.

b) Pour que $u:\deleted{\unsure{V}}I'\longrightarrow I$ soit une $W_0$-fibration, il f et s. que pour tte composante connexe $I'_\alpha$, de
\scanpage{186}\mspage{42}
% archive sheet number: 186

$I'$, posant $\alpha=u(\alpha')$, la flèche induite $I'_{\alpha'}\longrightarrow I_\alpha$ soit une $W_0$-équivalence universelle, \deleted{\unsure{i.e. soit}} (cf critère a) pour cette notion).

c) Pour que $u$ soit une $W_0$-fibration faible, il \uline{faut} que $\forall i\in I$, l'inclusion $\underset{\unsure{\hookrightarrow|u}}{I'}_i\hookrightarrow{}_{i\backslash}I'$ soit une $W_0$-équivalence i.e. $\pi_0(I'_i)\unsure{\simeq|\stackrel{\sim}{\longrightarrow}}\pi_0({}_{i\backslash}I')$ soit bijectif, \uline{et} que $i\longmapsto\pi_0({}_{i\backslash}I')$ soit un préfaisceau loc. constant sur $I$ (i.e. $\underline{\pi}_0(I'/I)$ soit un \uline{revêtement} étale de $I$). Cette condition est aussi suffisante, si $u$ est $W_0$-lisse. \deleted{\unsure{ici}} De plus (pour mémoire) $u$ est une $W_0$-fibration faible ssi $u^{\mathrm{opp}}$ l'est.

La plus grande différence entre la notion de $W_0$-fibration, et celle de $W_0$-fibration faible, me paraît que les morphismes $W_0$-parfaits (i.e. $W_0$-lisses et $W_0$-propres), et a fortiori les $W$-parfaits \deleted{$W$-propres} pour un local. fond. $W\subset W_0$, sont forcément des $W$-fibrations faibles, mais (\uline{\unsure{$=$|même} si $W=W_\infty$}, le plus fin de tous les loc. fondamentaux) \unsure{ce|ne} ne sont pas des $\unsure{W|W_0}$-fibrations. L'exemple le plus frappant à cet égard est, sans doute celui des \unsure{revêtements|revêtement} étales, qui sont $W_\infty$-parfaits, (donc
\scanpage{187}\mspage{43}
% archive sheet number: 187

des $W_\infty$-fibrations faibles et a fortiori des $W_0$-fibrations faibles), mais pas \unsure{de|des} $W_0$-fibrations ; car dans le cas où $I',I$ sont $0$-connexes, il faudrait que les fibres soient $0$-connexes \inserted{(cf b) ci-dessus)}, i.e. que $u$ soit un iso !

Une propriété \unsure{sympa.} des morphismes $W_0$-lisses (resp. $W_0$-propres), c'est qu'\deleted{\unsure{elles}} ce sont \inserted{des applications} \deleted{sont} \uline{ouvertes} \inserted{(resp. fermées)} pour les espaces topologiques sous-jacents), ce qui signifie en effet que les fibres locales i.e. des flèches localisées au but \inserted{$I'_{/\unsure{i'|u i'}}\longrightarrow I_{/\unsure{i|i'}}$} (resp colocales i.e. les flèches colocalisées au but $\,{}_{i'\backslash}I'\longrightarrow{}_{i\backslash}I$) sont surjectives, i.e. à fibres non-vides (ce qui est impliqué par la condition que celles-ci soient $0$-connexes).

Dans tout ça je n'ai \deleted{pas} encore guère parlé des flèches $W_0$-asphériques resp. co-asphériques, qui pourtant sont ici ma motivation \unsure{principale} pour introduire $W_0$ !

\uline{Proposition 3.2.} Conditions équivalentes sur la flèche $u:I'\longrightarrow I$ de $\mathrm{Cat}$ :

a) Pour tte $\mathcal{U}$-catégorie \inserted{$M$}, et tout foncteur $F:I\longrightarrow M$, la flèche canonique (sous réserve que les \deleted{\unsure{limites} \unsure{ind}} \inserted{petites} limites inductives sont \unsure{repr.|représ.} dans $M$)

\[
\underset{I'}{\underset{\longrightarrow}{\lim}}\,u^*\unsure{(}\deleted{\unsure{F}}F\unsure{)|,}\ \longrightarrow\ \underset{I}{\underset{\longrightarrow}{\lim}}\,F \tag*{\unsure{(3.2.1.)|(3.2,1.)}}
\]
\scanpage{188}\mspage{44}
% archive sheet number: 188

est iso.

$\deleted{\unsure{b}}a')$ Pour tte \deleted{\unsure{cat}} $\deleted{\unsure{M}}\mathcal{U}$-catégorie $M$ où les petites limites projectives sont représentables, \deleted{la flèche} \inserted{et tout $F\in M(I)=\underline{\mathrm{Hom}}(I^\circ,M)$}, \deleted{\unsure{pr} tout} \unsure{préfaisceau} sur $I$ à val. dans $M$, la flèche can\unsure{.}

\[
\underset{I}{\underset{\longleftarrow}{\lim}}\,F\longrightarrow\underset{I'}{\underset{\longleftarrow}{\lim}}\,u^*(\deleted{\unsure{G}}F)\tag*{(3.2.2.)}
\]

est iso.

$\deleted{\unsure{b}}a'')$ Comme a'), avec $M=(\mathrm{Ens})$ i.e. $F\in\mathrm{Ob}\,I^\wedge$.

% slanted
\margin{C'est une condition sur des morphismes de topos, savoir que $\forall F$ \deleted{dans} $\unsure{\mathscr{X}|X}$, $H^0(\mathscr{X},F)\overset{\sim}{\longrightarrow}H^0(\unsure{\mathscr{X}'},u^*F)$. Il y a notion similaire pour les $H^i$ $(i\leq n)$, \unsure{à dégager|à dég.} Notion des morphismes $W_0$-coasphériques des \unsure{catégories|objets} \unsure{fibrés|fléchés} d'un topos ! C'est \unsure{donc|dans} la \unsure{notion de dualité|notion de validité}, condition de \unsure{$g$|$G$}, \unsure{ou|vu} \unsure{comme} $W_0$-\unsure{asphérique}, qui \unsure{dérive|doute} de la notion \unsure{de} $\unsure{U^\wedge|U'}$ \unsure{induit} \unsure{des} signifie simplement que pour \unsure{tout} \unsure{objet} \unsure{dont le} $U^\wedge$, $\unsure{f^*|T^*}(U)$ est \unsure{$0$-connexe|connexe}. \unsure{Ici} \unsure{il convient de} \unsure{comparer} une notion \unsure{anciennement|``ancien''} de \deleted{\illegible{1}} $W_\infty$-asphéricité et \unsure{bien comprise}, \uline{\unsure{précisément}} \unsure{(ancien} $u$-asphérique), et (\deleted{\unsure{donc}} $W_\infty$-\unsure{coasphéricité}) \deleted{\unsure{est}} loc. \unsure{3.2.3.} \unsure{étant} celle évoquée plus haut \unsure{a''}.}

\margin{\uline{NB} On peut \unsure{intercaler|encore avoir} $a''')$ Comme a), mais avec $M=\mathrm{Ens}$ $a^{\mathrm{iv}})$ \deleted{\unsure{Comme}} a) avec le foncteur canonique $I'\longrightarrow\unsure{I'{}^\wedge|I^\wedge}$.}

b) La flèche $u$ est $W_0$-coasphérique, i.e. pour tout $i\in I$, $\,{}_{i\backslash}I'$ est $0$-connexe.

L'équivalence de a, a', a'' est $\pm$ tautologique, \unsure{compte tenu|vu} des définitions (cf \unsure{SGA 4 I 8.1}).

La a'') s'écrit, posant $\mathbb{D}=\mathrm{HOT}_{W_0}$, $\forall F\in\mathbb{D}(I)$, on veut que

\[
\begin{aligned}
H^{\unsure{\bullet|0|\cdot}}_{\mathbb D}(I,F)&\longrightarrow H^{\bullet}_{\mathbb D}(I',\deleted{\unsure{M}}u^*F)\\
\Vert\ (\mathrm{déf})&\qquad\Vert\ \mathrm{déf}\\
p_{I*}(F)&\qquad p_{I'*}(F)
\end{aligned}\tag*{(3.2.3.)}
\]

($p_I:I\longrightarrow e$ et $p_{I'}:I'\longrightarrow e$ les foncteurs canoniques) soit iso. On sait que cel\unsure{a} équivaut à $u$ $\mathbb{D}$-coasphérique. La démonstration est un jeu amusant entre $u_!$ \unsure{et} $u^*$, $u_*$ à la fin. Il faut exprimer que le $\mathrm{Hom}$ d'un objet $M$ de $\mathbb{D}(e)$ ($=\mathrm{Ens}$ en l'occurrence) vers la flèche (3.2.3) est bijective. Ici, il
\scanpage{189}\mspage{45}
% archive sheet number: 189

est plus naturel de prendre $M=e$ (ensemble ponctuel) et d'interpréter la flèche (3.2.3) comme

\[
\mathrm{Hom}(e_I,F)\longrightarrow\mathrm{Hom}(\underbrace{e_{I'}}_{\deleted{\unsure{*}}},u^*F)\tag*{(3.2.4)}
\]

et par adjonction \inserted{de $u_!,u^*$} le deuxième membre s'interprète comme $\mathrm{Hom}(u_!(e_{I'}),F)$, et la flèche est déduite de

\[
u_!e_{I'}=u_!u^*(e_I)\xrightarrow{\mathrm{adj.}}e_I \tag*{(\unsure{3.2.5})}
\]

en prenant le $\mathrm{Hom}$ dans $F$. Donc dire que (3.2.4) est iso, pour tt $F$, signifie que 3.2.5 est iso. \deleted{\unsure{Or les}} Passant aux valeurs des deux \unsure{mbs.|mbr.} en $i$, on trouve $\pi_0({}_{i\backslash}I')\overset{\sim}{\longrightarrow}1$, i.e. que $\,{}_{i\backslash}I'$ est $0$-connexe, \unsure{qed}.

% slanted
\margin{\uline{NB} Pour un énoncé \uline{complémentaire} de cofinalité cf \uline{XVIII} \unsure{prop.|par.} \unsure{4.9|4.}, 6.3 \unsure{2.} (p. 224--225, \unsure{238|228}--242) Il faudrait un énoncé qui coiffe les deux.}

\uline{Définition} 3.3. a) \deleted{\unsure{Pour que}} Quand les conditions équivalentes \inserted{de 3.2.} sont satisfaites, \inserted{\unsure{donc|i.e.} quand les $\,{}_{i\backslash}I'$ $(i\in\mathrm{Ob}\,I)$ sont $0$-conn.,} on dit que $u:\unsure{I''|I'}\longrightarrow I$ est \uline{cofinal}. On dit que $u$ est \unsure{cocofinal|cofinal} si $u^\circ:I'{}^\circ\longrightarrow I^\circ$ est cofinal, i.e. les localisées $I'_{/i}\unsure{/|,}$ de $I'$ sur $I$ $(i\in\mathrm{Ob}\,I)$ sont $0$-connexes.

b) Si $I'$ est une \uline{sous-catégorie} de $I$, on dit que $I'$ est \deleted{\unsure{cofinale dans}} $I$, si le foncteur d'inclusion est cofinal, \unsure{cocofinale} \deleted{\unsure{si le foncteur d'inclusion}} est \unsure{cocofinal|co-cofinal}.

\uline{Corollaire} 3.4. \textcircled{a} Pour que $u$ soit cofinal, \deleted{\unsure{donc}} il \uline{faut que} les cat. colocalisées $\,{}_{i\backslash}I'$ soient non vides, i.e. que l'on ait
\scanpage{190}\mspage{46}
% archive sheet number: 190

$\mathrm{F}\ 0)$ $\pi_0(\unsure{f|I|u})$ est bijectif

$\mathrm{F}\ 1)$ Pour tt $i\in\mathrm{Ob}\,I$, existe $i'\in I'$ et une flèche $i\longrightarrow u(i')$ de $I$.

% slanted
\margin{\uline{NB} F(1) implique déjà surjectivité \unsure{du|sur} $\pi_0$, donc sa bijectivité si $I'$ est $0$-connexe -- dans ce cas F0 \deleted{\unsure{est}} résulte de F1)}

% slanted
\margin{Si $I'$ est loc. \unsure{fblt} filtr. (= préfiltrante), alors pour la cofinalité de $u$, il \uline{suffit} qu'on ait F(0), F(1), F(2), et alors les $\,{}_{i\backslash}I'$ \unsure{sont} fblt filtrantes, et $I$ est préfiltrant tout comme $I'$ !}

b) Si $I'$ est fblt filtrante, pour que $u$ soit cofinal, il \uline{suffit} qu'elle satisfasse F1, ainsi que la condition :

$\mathrm{F}\ 2)$ Pour toute double flèche

\diagram{\begin{tikzcd}[ampersand replacement=\&]
i \arrow[r,shift left=.6ex,"\beta"] \arrow[r,shift right=.6ex,"\gamma"'] \& u(i')
\end{tikzcd}}

\deleted{existe} (avec $i\in\mathrm{Ob}\,I$, $i'\in\mathrm{Ob}\,I$) existe une flèche $\deleted{\unsure{b'}}i'\xrightarrow{\alpha'}j'$ de $I'$, de telle façon que $u(\alpha')$ égalise $\beta,\gamma$, i.e. $u(\alpha')\beta=u(\alpha')\gamma$.

Et ces conditions F1, F2 impliquent que les $\,{}_{i\backslash}I'$ sont fblt filtrantes, et $I$ aussi.

% slanted
\margin{Si $I'$ est loc filtr.-te il est cofinal ssi \deleted{\unsure{elle}} satisfait \deleted{\unsure{alors}} F0, F1, F2, et \deleted{\unsure{alors}} les $\,{}_{i\backslash}I'$ sont filtr.-tes, et alors $I$ est loc. filtrante.}

c) Si $I'$ est filtrante, alors $u$ est cofinal \uuline{ssi} il satisfait les deux conditions F1, F2, ou \deleted{\unsure{est}} encore ssi les $\,{}_{i\backslash}I'$ sont \deleted{\unsure{fblt}} \uline{filtrantes} (et non seulement $0$-connexes). Et cela implique que $\unsure{I|I'}$ est \inserted{également} filtrante.

d) Si $I$ est \inserted{\unsure{pré}}filtrante \inserted{(= loc. \unsure{fblt.|fbt.} filtrante, i.e. PS.1)} \inserted{et $u$} \uline{pl. fidèle}, pour que $u$ soit cofinal, il \uline{f et s} que $u$ satisfasse la condition F1 plus haut. Et alors $I'$ est également \deleted{filtrante,} \deleted{loc. filtrante} préfiltrante, et les $\,{}_{i\backslash}I'$ sont fblt. filtrantes. Si de plus $I$ est \deleted{connexe, i.e. fblt filtrante} \deleted{\unsure{(resp.} si $I$ est} filtrante\deleted{\unsure{s}}, \inserted{les $\,{}_{i\backslash}I'$ et} $I'$ \unsure{l'est|l'ét} également.

(Comparer SGA 4 I 8.1.3.)
\scanpage{191}\mspage{47}
% archive sheet number: 191

Dém. : a) \inserted{F1)} est évident, car $0$-connexe $\Rightarrow$ non vide. \inserted{F0) car $W_0$-coasph. $\Rightarrow$ $W_0$-équiv. i.e. $\pi_0(u)$ bij.}

Pour \uline{prouver b)}, il suffit de prouver que F1, F2 (si $I'$ est fblt filtrante) implique que les $\,{}_{i\backslash}I'$ sont fblt filtrantes (a fortiori $0$-connexes) \inserted{et de plus \deleted{\unsure{le fait}} que $I$ \unsure{l'est} également}. On sait déjà qu'ils sont non vides par F1, reste à voir que si $(i',\alpha),(j',\beta)$ sont dans $\,{}_{i\backslash}I'$, alors ils sont ``majorés'' par un $3^{\mathrm{o}}$ objet de $\,{}_{i\backslash}I'$.

\diagram{\begin{tikzcd}[ampersand replacement=\&,row sep=1.3em,column sep=1.5em]
\& i'\arrow[rr,"r"]\arrow[dd,no head] \& \& k'\arrow[r,dashed,"t"]\arrow[dd] \& l'\arrow[dd,no head]\\
\& \& j'\arrow[ur,"s"]\arrow[dd,no head,crossing over] \& \&\\
\& u(i')\arrow[rr,"{u(r)}"] \& \& u(k')\arrow[r,dashed,"{u(t)}"] \& \deleted{\unsure{u(l')u(l')|u(l'),u(l')}}\\
i\arrow[ur,"\alpha"]\arrow[rr,"\beta"'] \& \& u(j')\arrow[ur,"{u(s)}"'] \& \&
\end{tikzcd}}

Comme $I'$ est fblt filtr., on trouve

\diagram{\begin{tikzcd}[ampersand replacement=\&,row sep=small]
i'\arrow[rr,"r"] \& \& k'\\
\& j'\arrow[ur,"s"'] \&
\end{tikzcd}$,$}

\deleted{\unsure{d'où}} \deleted{\unsure{les deux flèches}} \inserted{$u(r),u(s)$} de $I$ (cf diagramme), appliquons F2 : la double-flèche $u(r)\alpha,u(s)\beta:i\rightrightarrows u(k')$, on trouve $k'\xrightarrow{t}l'$ (en pointillés sur diagramme) telle que $(u(t)u(r))\alpha=(u(t)u(s))\beta$, \inserted{soit $\gamma$,} ce qui montre que l'objet $(l',\gamma:\unsure{i'|i}\longrightarrow u(l'))$ majore $(i',\alpha)$ et $(j',\beta)$.

\inserted{Si $I'$ est fblt filtrant et $u$ satisfait F1, alors $I$ est fblt filtrant, c'est immédiat.}

\uuline{Prouvons c).} Il \unsure{reste} à prouver que si $I'$ est \uline{filtrante} (non seulement fblt filtrante), \deleted{alors} \deleted{\unsure{al}} \inserted{si} $u$ cofinal, alors F2 est vérifiée, et \deleted{que $I$} \inserted{de plus} \deleted{et plus} $\,{}_{i\backslash}I'$ \inserted{et $I$} sont filtrantes (pas seulement fblt. filtrantes). \deleted{\unsure{Or}} Mais notons \deleted{que} tout de suite que $\,{}_{i\backslash}I'\longrightarrow I'$ est coétale, donc $I'$ filtrante \inserted{donc loc. filtr.-te} implique que $\,{}_{i\backslash}I'$ est aussi loc. filtrante ; comme par hypothèse les $\,{}_{i\backslash}I'$ sont $0$-connexes, elles sont donc filtr.-tes. La condition F2 \deleted{\unsure{soit}} \unsure{sort|suit} facilement à présent :
\scanpage{192}\mspage{48}
% archive sheet number: 192

\diagram{\begin{tikzcd}[ampersand replacement=\&,row sep=1.25em,column sep=1.15em]
\& \& i'\arrow[dr,"{\deleted{\unsure{s}}\,\unsure{r}}",Rightarrow]\arrow[dd] \& \&\\
\& \unsure{i''|i'}\arrow[rr,Rightarrow,"s"]\arrow[dd,no head] \& \& j'\arrow[r,"t"]\arrow[r,dotted,line width=1.1pt]\arrow[dd,no head] \& k'\arrow[dd,no head]\\
i\arrow[rr,"\alpha"]\arrow[dr,"\beta"'] \& \& u(i')\arrow[dr,"{u(r)}"] \& \&\\
\& u(\unsure{i'|i''})\arrow[rr,"{u(s)}"]\arrow[dr,shorten >=1.5em] \& \& u(j')\arrow[r,dashed,"{u(t)}"] \& u(k')\\
\& \& {} \& \& 
\end{tikzcd}}

Comme $\,{}_{i\backslash}I'$ est filtrant donc fblt filtrant \inserted{(PS1)}, on peut trouver

\diagram{\begin{tikzcd}[ampersand replacement=\&]
i'\arrow[r,shift left=.6ex,no head,"r"]\arrow[r,shift right=.6ex,no head,"s"'] \& j'
\end{tikzcd}}

telles que

\[
u(r)\alpha=u(s)\beta
\]

et comme $I'$ est filtrant on peut trouver $t:j'\longrightarrow k'$ qui égalise la double-flèche $(r,s)$ :

\[
tr=ts\quad(\overset{\mathrm{def}}{=}\tau),
\]

alors

\[
\underbrace{u(t)u(r)}_{u(tr)}\alpha=\underbrace{u(t)u(s)}_{u(ts)}\beta\qquad\text{i.e.}\quad u(\tau)\alpha=u(\tau)\beta,
\]

i.e. on a vérifié F2.

Reste à prouver que $I$ est filtrante (si $I'$ l'est et $u:I'\unsure{\longrightarrow|\hookrightarrow}I$ est cofinal). On sait déjà que $I$ est fblt filtrante par b), reste à vérifier la condition de la double flèche

\diagram{\begin{tikzcd}[ampersand replacement=\&]
i\arrow[r,shift left=.6ex,"\alpha"]\arrow[r,shift right=.6ex,"\beta"'] \& j
\end{tikzcd}}

(PS2). Soit $\unsure{j'|i'}$ dans $I'$ et $\unsure{\gamma}:\!j\longrightarrow\deleted{\unsure{j}}u(j')$, \deleted{alors $(j',\unsure{\gamma}\alpha)$ et $i'$} et considérons

\diagram{\begin{tikzcd}[ampersand replacement=\&]
i\arrow[r,shift left=.6ex,"{\deleted{\unsure{\beta}}\gamma\alpha}"]\arrow[r,shift right=.6ex,"{\gamma\beta}"'] \& u(\unsure{i'|j'})
\end{tikzcd}$,$}

par hyp. F2 $\exists\,t':i'\longrightarrow j'$ dans $I'$ \unsure{telle|telles} que
\scanpage{193}\mspage{49}
% archive sheet number: 193

$u(\unsure{r|t'})(\gamma\unsure{\alpha'|\alpha})=u(r)(\gamma\beta)$, donc $u(\unsure{r})\gamma$ égalise $(\alpha,\beta)$, qed

La vérification de d) \unsure{en|un} \unsure{AQT}.

4. \uline{Foncteurs biasphériques}

Soit encore

\[
u:I'\longrightarrow I
\]

une flèche de $\mathrm{Cat}$, soit $M$ une catégorie, d'où \deleted{où les lim} \deleted{supposons que}

\[
u^*:\ M(I)\longrightarrow M(I')\ ;\tag*{(4.1.)}
\]

\deleted{\unsure{pour qu'}} nous cherchons des conditions \inserted{sur $u,M$} pour que ce foncteur soit pleinement fidèle. Si le foncteur adjoint à gauche $u_!$ (resp. adjoint à droite $u_*$) existe, \unsure{On|on} sait qu'il revient au $=$ de dire que le morphisme d'adjonction

\[
u_!u^*F\longrightarrow F\tag*{(4.2)}
\]

resp

\[
u_*u^*F\longleftarrow F\tag*{(4.\unsure{2|1} bis)}
\]

soit iso pour tout $F$ dans $M(I)$. Disons alors que $u\deleted{\unsure{f}}$ est $M$-connexe, ou $M$-biasphérique, cette terminologie sous-entendant que c'est là une propriété qui se réfère au dérivateur $\mathbb{D}_M=S\longrightarrow M(S)$ sur $\mathrm{Cat}$ défini par $M$. Notons que

\[
\begin{aligned}
&\text{Si }u\text{ est }M\text{-}\deleted{\unsure{\mathrm{bi}}}\text{connexe, alors il est aussi}\\
&\deleted{\unsure{Hom}} M(S)\text{-connexe, \unsure{quel|ceci} pour }\deleted{\text{tout}} S\text{ dans }\mathrm{Cat},
\end{aligned}\tag*{(4.3.)}
\]

puisque $M(S)(I)\simeq\deleted{\unsure{H}}M(I)\deleted{\unsure{\}}}(S)\unsure{=|\simeq}\underline{\mathrm{Hom}}(S^\circ,M(I))$,
\scanpage{194}\mspage{50}
% archive sheet number: 194

et

\[
\begin{aligned}
&\text{Si }M'\hookrightarrow M^{\unsure{\deleted{\prime}|E|F}}\text{ est pl. fid, alors }M^{\unsure{\deleted{\prime}|E|\leftarrow}}\text{-connexe}\\
&\text{implique }M'\text{-connexe,}
\end{aligned}\tag*{(4.4.)}
\]

enfin, inversement

\[
\begin{aligned}
&\text{pour que }u\text{ soit }M\text{-connexe, il f. et s.}\\
&\text{que pour tte \uline{petite} sous-catégorie}\\
&\text{pleine }M'\text{ de }M,\ u\text{ soit }M'\text{ connexe.}
\end{aligned}\tag*{(4.5.)}
\]

Comme toute petite $M'$ se plonge dans $M'{}^\wedge\deleted{\unsure{,}}=\deleted{\unsure{M}}\mathrm{Ens}(M')$, il résulte de \unsure{.} (4.3)(4.4)(4.5) que \deleted{pour \unsure{tout}}

\[
\begin{aligned}
&\text{Si }u\text{ est Ens-connexe, il est }M\text{-connexe}\\
&\text{pour tte catégorie }M,
\end{aligned}\tag*{(4.6)}
\]

car c'est le cas si $M$ est petite par (4.3), donc pour tte $M$ par (4.5).

Disons qu'un foncteur $u$ \inserted{dans $\mathrm{Cat}$} est $M$-\uline{cofinal} \inserted{ou $M$-coasphérique} si \unsure{le|la} \deleted{\unsure{foncteur|fonction}} \unsure{hom|homo} fonctoriel

\[
\underset{\deleted{\unsure{*}}I}{\underset{\longleftarrow}{\lim}}\,F\longrightarrow
\underset{I'}{\underset{\longleftarrow}{\lim}}\,u^*(F)\tag*{(4.7)}
\]

(flèche dans $\underline{\mathrm{Hom}}(M(I),M)$ si les $\underset{\longleftarrow}{\lim}$ de types $I,I'$ existent dans $M$, sinon flèche dans $\underline{\mathrm{Hom}}(M(I),M^\wedge)$) est \deleted{\unsure{un}} iso. En se ramenant au cas de petites catégories $M$, \inserted{\unsure{chacune|disons|dûment}} plongées dans $M^\wedge$, on voit par le critère (4.2.bis) que

\[
u\text{ est }M\text{-connexe si pour tt }i\in\mathrm{Ob}\,I,\tag*{(4.8.)}
\]
\scanpage{195}\mspage{51}
% archive sheet number: 195

\deleted{\unsure{les}} \inserted{le} foncteur induit $I'_{/i}\longrightarrow I_{/i}$ est $M$-cofinal.

La réciproque n'est pas claire, car \deleted{la condition} \inserted{$u$ est $M$-connexe}, \deleted{dite exprime que} \deleted{\unsure{exprime} seulement que, pour} si et seulement si pour tout $i\in\mathrm{Ob}\,I$ et tout \inserted{$M$-}préfaisceau $F_i$ sur $I_{/i}$ \deleted{provenant de} \uline{induit} \deleted{des \unsure{images} \unsure{essentielles}} \uline{\unsure{par|de} un préfaisceau \inserted{$F$} sur} $I$, l'hom

\[
\underset{I_{/i}}{\underset{\longleftarrow}{\lim}}F_i\longrightarrow\underset{I'_{/i}}{\underset{\longleftarrow}{\lim}}(u_{/i})^*(F_i)
\]

est iso. On ne peut donc affirmer d'emblée que cela implique itou pour \uline{tout} $F_i$ sur $I_{/i}$, que si on sait que $(\unsure{U|u}_{/i})^{\unsure{*|\wedge}}:M(I)\longrightarrow M(I_{/i})$ est ess. surjectif. Il en est ainsi pour tout $i$ quand $I$ est \uline{ordonné}, \deleted{mais} car alors $I_{/i}\longrightarrow I$ est pl. fidèle, donc $M(I_{/i})\longrightarrow M(I)$ pl. fidèle, donc $M(I)\longrightarrow M(I_{/i})$ un foncteur de localisation, \deleted{Par contre, ce n'est pas le cas} en particulier, ess. surjectif.

\deleted{\deleted{\unsure{pour}} En fait, il n'y a pas à supposer forcément (pour la réciproque à (4.8.) que les \inserted{$\alpha_i^*:$} $M(I)\longrightarrow M(I_{/i})$ soient ess. surjectifs, mais (sous réserve d'existence de $\alpha_{i*}$) que \uline{$\alpha_{i*}$ soit conservatif}. Car \deleted{\unsure{pour}} il faut exprimer que \deleted{$\unsure{u_*u^*F}\unsure{\longrightarrow|\longleftarrow}F\unsure{(i)}$} \inserted{$\Gamma(I'_{/i},\unsure{\alpha^*|u^*}F)$} $\longleftarrow\Gamma(I_{/i},F)$ est iso, pour tt $F$ dans $M(I)$, or ceci}
\scanpage{196}\mspage{52}
% archive sheet number: 196

\deleted{s'interprète comme}

Mais quand $M=(\mathrm{Ens})$, cette perplexité disparaît. On trouve la

\uline{Proposition} 4.1. Soit $u:I'\to I$ une flèche de $\mathrm{Cat}$. Conditions équivalentes

a) Pour tte catégorie $M$, $u^*_M:M(I)\to M(I')$ est pl. fid. (i.e. $u$ est $M$-connexe)

a') Pour tte catégorie $M$\unsure{,} où les limites inductives de type $I,I'$ existent, \unsure{la|les} \unsure{flèche|flèches} d'adjonction pour $F$ dans $M(I)$

\[
u_!u^*F\xrightarrow{\sim}\unsure{F|G}
\]

est iso.

a'') Dual de a'), avec flèche d'adjonction

\[
u_*u^*\deleted{\unsure{F|G}}F\xleftarrow{\sim} F.
\]

b) b') b'') Comme a) a') a''), avec $M=(\mathrm{Ens})$

c) Comme b', avec $F$ \deleted{\unsure{des|de}} représentable, i.e de la forme $\deleted{\unsure{i|j}}j\longmapsto\operatorname{Hom}_I(j,i)$.

\deleted{c') Comme b'', avec $F$ représentable}

d) Pour tout $i$ dans $I$, $I'_{/i}\to I_{/i}$ est cofinal

d') Pour tout $i$ dans $I$, $\,{}_{i\backslash}I'\to{}_{i\backslash}I$ est cocofinal.

d'') Pour toute flèche $i_0\xrightarrow{\alpha}i_1$ de $I$, la catégorie $\,{}_{i_0\backslash}\underset{\alpha}{I'}_{/i_1}\overset{\mathrm{def}}{=}I'\times_I\underbrace{{}_{i_0\backslash}\underset{\alpha}{I}_{/i_1}}_{\substack{\text{catégorie biinduite}\\(i_0,\alpha)\backslash(I_{/i_1})\,\simeq\,({}_{i_0\backslash}I)/(i_1,\alpha)}}$ est $0$-connexe.
